\documentclass[11pt,a4paper]{article}

%% PACHETE MATEMATICE
\usepackage{amsmath,amssymb,amsfonts}
\usepackage{amsthm}
\usepackage{mathpazo}
\usepackage{eulervm}
\usepackage{enumerate}
\usepackage{color}
\usepackage{graphicx}
\usepackage[all]{xy}
\usepackage{xcolor}
\usepackage{framed}
\usepackage[unicode]{hyperref}
\setcounter{MaxMatrixCols}{20}
\usepackage{multicol}
% \usepackage[romanian]{babel}


%% ASPECT
\linespread{1.05}
\usepackage[T1]{fontenc}
\usepackage[utf8x]{inputenc}
\usepackage[margin=2cm]{geometry}
% \usepackage[color]{showkeys}
\usepackage{anyfontsize}
\usepackage{titlesec}
\titleformat*{\section}{\large\bfseries}

\usepackage{fancyhdr}
\pagestyle{fancy}
\rhead{\small \color{gray} Notițe de Adrian Manea}
\lhead{\small \color{gray} ETTI-AN1, 2017---2018, C. Ghiu}


\usepackage[autostyle = false, style = german]{csquotes}
\usepackage[romanian]{babel}
\newcommand\qq{\enquote}

%% COMENZILE MELE
\renewcommand*{\proofname}{Dem.:}
\newcommand\bb{\mathbb}
\newcommand\dr{\mathrm}
\newcommand\kal{\mathcal}
\newcommand\fk{\mathfrak}
\newcommand\op{\oplus}
\newcommand\ot{\otimes}
\newcommand\ds{\displaystyle}
\newcommand\id{\indent}
\newcommand\nid{\noindent}
\newcommand\seq{\subseteq}
\newcommand\sneq{\subsetneq}
\newcommand\speq{\supseteq}
\newcommand\spneq{\supsetneq}
\newcommand\rar{\rightarrow}
\newcommand\Rar{\Rightarrow}
\newcommand\xrar{\xrightarrow}
\newcommand\lar{\leftarrow}
\newcommand\Lar{\Leftarrow}
\newcommand\xlar{\xleftarrow}
\newcommand\lrar{\leftrightarrow}
\newcommand\Lrar{\Leftrightarrow}
\newcommand\ideal{\trianglelefteq}
\newcommand\ai{\text{ a.\^{i}. }}
\newcommand\sk{(\textit{Schi\c{t}\u{a}}) }
\newcommand\rhup{\rightharpoonup}
\newcommand\rhdn{\rightharpoondown}
\newcommand\lhup{\leftharpoonup}
\newcommand\lhdn{\leftharpoondown}
\newcommand\vid{\emptyset}
\newcommand\wh{\widehat}
\newcommand\ol{\overline}
\newcommand\NN{\mathbb{N}}
\newcommand\RR{\mathbb{R}}
\newcommand\ZZ{\mathbb{Z}}
\newcommand\QQ{\mathbb{Q}}
\newcommand\CC{\mathbb{C}}

%% STILURILE MELE
\newtheoremstyle{dotless}{}{}{\itshape}{}{\bfseries}{:}{ }{}

\theoremstyle{dotless}
\newtheorem{theorem}{Teorem\u{a}}[section]
\newtheorem{proposition}{Propozi\c{t}ie}[section]
\newtheorem{lemma}{Lem\u{a}}[section]
\newtheorem{corollary}{Corolar}[section]
\newtheorem{exercise}{Exerci\c{t}iu}

\newtheoremstyle{dotlessdr}{}{}{}{}{\bfseries}{:}{ }{}
\theoremstyle{dotlessdr}
\newtheorem{example}{Exemplu}[section]
\newtheorem{definition}{Defini\c{t}ie}[section]
\newtheorem{remark}{Observa\c{t}ie}[section]




\begin{document}
% \definecolor{labelkey}{RGB}{222,0,0}

\begin{center}
  {\large\textbf{Seminar 4 \\ Serii Fourier și recapitulare}}
\end{center}

\section{Serii Fourier}
\label{sec:fourier}

Pentru dezvoltarea în serie Fourier (care se poate aplica atunci cînd seriile
Taylor sînt imposibile), trebuie satisfăcute \emph{condițiile Dirichlet}:
\begin{enumerate}[(D1)]
\item Funcția trebuie să fie periodică;
\item Funcția trebuie să aibă un număr finit de discontinuități;
\item Funcția trebuie să aibă un număr finit de extreme într-o perioadă;
\item $\ds\int_0^T |f(x)| dx$ trebuie să fie convergentă.
\end{enumerate}

Convențional, seria se scrie:
\[
  f(x) = \frac{a_0}{2} + \sum_{r = 1}^\infty \Bigg[ a_r \cos \Big( \frac{2\pi rx}{L} \Big) +
  b_r \sin \Big( \frac{2\pi rx}{L} \Big) \Bigg],
\]
unde $a_r, b_r, a_0$ sînt \emph{coeficienții Fourier}, iar $L$ este perioada.

Coeficienții se calculează:
\begin{align*}
  a_r &= \frac{2}{L} \int_{x_0}^{x_0 + L} f(x) \cos \Big( \frac{2\pi rx}{L} \Big) dx \\
  b_r &= \frac{2}{L} \int_{x_0}^{x_0 + L} f(x) \sin \Big( \frac{2\pi rx}{L} \Big) dx.
\end{align*}

De obicei, se ia $x_0 = 0$ sau $x_0 = -L/2$.

\begin{theorem}[Parseval]\label{thm:parseval}
  \[
    \frac{1}{L} \int_{x_0}^{x_0 + L} |f(x)|^2 dx =
    \sum_{r = -\infty}^\infty |c_r|^2 =
    \Big( \frac{a_0}{2} \Big)^2 + \frac{1}{2} \sum_{r = 1}^\infty (a_r^2 + b_r^2).
  \]
\end{theorem}

De exemplu, putem calcula $\sum r^{-4}$.

Luăm funcția $f(x) = x^2$ și calculăm media funcției $f^2(x)$ pe $-2 < x \leq 2$:
\[
  \frac{1}{4} \int_{-2}^2 x^4 dx = \frac{16}{5}.
\]
Acum calculăm membrul drept din Parseval:
\[
  \Big( \frac{a_0}{2} \Big)^2 + \frac{1}{2} \sum_{r = 1}^\infty (a_r^2 + b_r^2) =
  \Big( \frac{4}{3} \Big)^2 + \frac{1}{2} \sum_{r \geq 1} \frac{16^2}{\pi^4 r^4}.
\]
Egalăm cele două expresii și obținem $\sum_{r \geq 1} \frac{1}{r^4} = \frac{\pi^4}{90}$.


Folosind forma polară a unui număr complex și formula lui Euler, putem scrie seria:
\begin{align*}
  \widehat{f} &= \frac{a_0}{2} + \sum_{n = 1}^\infty (a_n \cos nt + b_n \sin nt), \\
  a_n &= \frac{1}{\pi} \int_0^{2\pi} f(t) \cos nt dt, n \geq 0 \\
  b_n &= \frac{1}{\pi} \int_0^{2\pi} f(t) \sin nt dt, n \geq 1.
\end{align*}

\begin{lemma}[Riemann]\label{le:riemann-fourier}
  Dacă $f$ este integrabilă, atunci:
  \[
    \lim_{n \to \infty} a_n = \lim_{n \to \infty} b_n = 0.
  \]
\end{lemma}

\begin{theorem}[Dirichlet]\label{thm:dirichlet-fourier}
  Dacă $f : \RR \to \RR$ este o funcție periodică, de perioadă $2\pi$, măsurabilă,
  mărginită, avînd cel mult un număr finit de discontinuități de speța întîi și cu
  derivate laterale în orice punct, atunci seria Fourier converge în fiecare punct
  $x \in \RR$ la:
  \[
    \frac{1}{2} \big( f(x + 0) + f(x - 0) \big).
  \]
  În particular, dacă $f$ este chiar continuă, are loc:
  \[
    f(t) = \frac{a_0}{2} + \sum_{n = 1}^\infty (a_n \cos nt + b_n \sin nt).
  \]
\end{theorem}



\subsection{Exerciții}

Să se dezvolte în serie Fourier:

1. $f(x) = x$, pe $(-\pi, \pi)$.

\emph{Soluție:} Funcția este impară, deci $a_k = 0, \forall k \geq 0$, iar:
\[
  b_k = \frac{2}{\pi} \int_0^\pi x \sin kx dx =
  \frac{2 \cdot (-1)^{k+1}}{k}, k > 0.
\]
Egalitatea se obține ținînd cont că $\sin k \pi = 0, \cos k\pi = (-1)^k$.

Deci, pentru orice $x \in \RR$, avem:
\[
  x = 2 \sum_{k = 1}^\infty \frac{(-1)^{k + 1}}{k} \sin kx.
\]
Pentru $x = \frac{\pi}{2}$, avem:
\[
  \frac{\pi}{4} = 1 - \frac{1}{3} + \frac{1}{5} - \frac{1}{7} + \dots.
\]

\vspace{1cm}

2. $f(x) = \pi^2 - x^2$, pe $(-\pi, \pi)$.

\emph{Soluție:} Funcția este pară, deci $b_k = 0, \forall k > 0$ și:
\begin{align*}
  a_0 &= \frac{2}{\pi} \int_0^{\pi} (\pi^2 - x^2) dx = \frac{4\pi^2}{3} \\
  a_k &= \frac{2}{\pi} \int_0^{\pi} (\pi^2 - x^2) \cos kx dx \\
      &= \frac{4 \cdot (-1)^{k - 1}}{k^2}, k > 0.
\end{align*}

Deci:
\[
  \pi^2 - x^2 = \frac{2\pi^2}{3} + 4 \cdot \sum_{k = 1}^\infty \frac{(-1)^{k-1}}{k^2} \cos kx.
\]
Sînt îndeplinite condițiile din teorema Dirichlet, deci descompunerea este
valabilă pentru $x \in [-\pi, \pi]$.

În particular, dacă $x = \pi$, avem formula:
\[
  \sum_{k = 1}^\infty \frac{1}{k^2} = \frac{\pi^2}{6}.
\]

\vspace{1cm}

3.
\[
  f(x) =
  \begin{cases}
    ax, & x \in (-\pi, 0) \\
    bx, & x \in [0, \pi).
  \end{cases}
\]

Coeficienții:
\begin{align*}
  a_0 &= \frac{1}{\pi} \int_{-\pi}^\pi f(x) dx = \frac{\pi}{2} (b - a) \\
  a_n &= \frac{1}{\pi} \int_{-\pi}^\pi f(x) \cos nx dx \\
      &= \frac{1}{\pi} \Big( \int_{-\pi}^0 ax \cos nx dx + \int_0^\pi bx \cos nx dx \Big) \\
  b_n &= \frac{1}{\pi} \int_{-\pi}^\pi f(x) \sin nx dx \\
      &= \frac{1}{\pi} \Big( \int_{-\pi}^0 ax \sin nx dx + \frac{1}{\pi} \int_0^x ax \sin nx dx \Big)
\end{align*}

Rezultă, prin calcule (integrare prin părți):
\[
  a_k = \frac{a - b}{\pi} \cdot \frac{1 - (-1)^k}{k^2}, \quad
  b_k = (a + b) \cdot \frac{(-1)^{k - 1}}{k}.
\]

\textbf{Observație:} Putem folosi această dezvoltare pentru a calcula suma
seriei numerice $\sum \frac{1}{(2n - 1)^2}$, care este convergentă.

Funcția $f$ este continuă pentru $t = 0$ și, din Dirichlet, suma seriei Fourier
în punctul $t = 0$ este egală cu $f(0)$. Obținem:
\[
  0 = f(0) = - \frac{a - b}{4} \pi + \frac{2(a - b)}{\pi}
  \sum_{n = 1}^\infty \frac{1}{(2n - 1)^2}.
\]

Rezultă $\sum_{n = 1}^\infty \frac{1}{(2n - 1)^2} = \frac{\pi^2}{8}$.

\vspace{1cm}

4. $f(x) = e^{ax}, a \neq 0$, pe $(-\pi, \pi)$.

\emph{Soluție:} Coeficienții:
\begin{align*}
  a_0 &= \frac{1}{\pi} \int_{-\pi}^\pi e^{ax} dx = \frac{2}{a\pi} \sinh a\pi \\
  a_n &= \frac{1}{\pi} \int_{-\pi}^\pi e^{ax} \cos nx dx \\
      &= (-1)^n \frac{1}{\pi} \cdot \frac{2a}{a^2 + n^2} \sinh a\pi \\
  b_n &= \frac{1}{\pi} \int_{-\pi}^\pi e^{ax} \sin nx dx \\
      &= (-1)^{n - 1} \frac{1}{\pi} \cdot \frac{2n}{a^2 + n^2} \sinh a\pi.
\end{align*}

Rezultă:
\[
  e^{ax} = \frac{2\pi} \sinh a\pi \cdot
  \Bigg[ \frac{1}{2a} + \sum_{n = 1}^\infty \frac{(-1)^n}{a^2 + n^2} \cdot
  (a \cos nx - n \sin nx) \Bigg].
\]

\vspace{1cm}

5. $f(x) = |x|$, pe $[-\pi, \pi]$.

\emph{Soluție:} Funcția este pară, deci $b_n = 0, \forall n \geq 1$.
\begin{align*}
  a_0 &= \frac{2}{\pi} \int_0^\pi x dx = \pi \\
  a_n &= \frac{2}{\pi} \int_0^\pi x \cos nx dx =
        \frac{2}{\pi n^2} [(-1)^n - 1].
\end{align*}

Rezultă:
\[
  |x| = \frac{\pi}{2} - \frac{4}{\pi} \cdot
  \sum_{n = 1}^\infty \frac{\cos(2n - 1) x}{(2n - 1)^2}.
\]

\vspace{1cm}

6. Să se dezvolte în serie de sinusuri funcția $f(x) = \frac{|x|}{x}$, definită
în intervalul $(0, 1)$.

\emph{Soluție:} Prelungim funcția impar față de origine:
\[
  f(x) =
  \begin{cases}
    1,  & 0 < 1 < 1\\
    0,  & x = 0 \\
    -1, & -1 < x < 0
  \end{cases}
\]
Calculăm coeficienții Fourier ai acestei funcții periodice impare, definite
pe $(-1, 1)$:
\begin{align*}
  a_n &= 0, \forall n \geq 0 \\
  b_n &= \frac{2}{1} \int_0^1 1 \cdot \sin \frac{n \pi x}{1} dx \\
      &= \frac{2}{\pi} \cdot \frac{1 - (-1)^n}{n} \\
      &= \begin{cases}
        \frac{4}{\pi} \cdot \frac{1}{2n +1}  & n = 2k + 1 \\
        0, & n = 2k
      \end{cases}
\end{align*}

Rezultă:
\[
  f(x) = \frac{4}{\pi} \sum_{n = 0}^\infty \frac{1}{2n + 1} \sin(2n + 1) \frac{\pi x}{1}.
\]

\textbf{Observație:} 1 poate fi înlocuit cu $l$, orice perioadă aleasă.

\vspace{1cm}

7. Să se demonstreze formula:
\[
  \frac{\pi - x}{2} = \sum_{n \geq 1} \frac{\sin nx}{n}, \forall x \in (0, 2\pi).
\]

\emph{Soluție:} Considerăm funcția $f(x) = \frac{\pi - x}{2}, x \in [0, 2\pi]$,
prelungită prin periodicitate la $\RR$. Calculăm coeficienții Fourier:
\begin{align*}
  a_0 &= \frac{1}{\pi} \int_0^{2 \pi} \frac{\pi - x}{2} dx \\
      &= \frac{1}{2\pi} \Big( \pi x - \frac{x^2}{2} \Big) \Big|_0^{2 \pi} = 0; \\
  a_n &= \frac{1}{\pi} \int_0^{2\pi} \frac{\pi - x}{2} \cos nx dx \\
      &= \frac{\pi - x}{\sin nx} 2n\pi \Big|_0^{2 \pi} -
        \frac{1}{2n \pi} \int_0^{2\pi} \sin nx dx = 0, \forall n \geq 1 \\
  b_n &= \frac{1}{n} \int_0^{2\pi} \frac{\pi - x}{2} \sin nx dx \\
      &= \frac{-(\pi - x) \cos nx}{2n \pi} \Big|_0^{2\pi} -
        \frac{1}{2n\pi} \int_0^{2\pi} \cos nx dx \\
      &= \frac{1}{n}, \forall n \geq 1.
\end{align*}

Aplicăm acum teorema lui Dirichlet și obținem:
\[
  \frac{\pi - x}{2} = \sum_{n \geq 1} \frac{\sin nx}{n}, \forall x \in (0, 2\pi).
\]
Pentru $x = 0, 2\pi$, funcția nu este continuă. Seria trigonometrică
asociată are suma 0.


%%%%%%%%%%%%%%%%%%%%%%%%%%%%%%%%%%%%%%%%%%%%%%%%%%%%%%%%%%%%%%%%%%%%%%

\section{Exerciții recapitulative}
\label{sec:recap}

\textbf{Serii numerice}:

1. Decideți convergența următoarelor serii cu termenul general dat de:
\begin{enumerate}[(a)]
\item $x_n = \frac{n!}{n^{2n}}$ \hfill (C, raport);
\item $x_n = \Big( 1 - \frac{3 \ln n}{2n} \Big)^n$ \hfill(C, logaritmic);
\item $x_n = \Big( \frac{1}{\ln n} \Big)^{\ln(\ln n)}$ \hfill(D, logaritmic);
\item $x_n = \frac{1}{n \ln^2 n}$ \hfill (C, integral);
\item $x_n = \frac{1}{n \sqrt[n]{n}}$ \hfill (D, comparație $\sim \frac{1}{n}$);
\item $x_n = a^{\ln n}, a > 0$ \hfill (discuție, Raabe);
\item $x_n = \frac{(-1)^n \sqrt{n} + 1}{n}$ \hfill (D, Leibniz + armonică);
\item $x_n = n^2 e^{-\sqrt{n}}$ \hfill (C, logaritmic).
\end{enumerate}

\vspace{1cm}
%%%%%%%%%%%%%%%%%%%%%%%%%%%%%%%%%%%%%%%%%%%%%%%%%%%%%%%%%%%%%%%%%%%%%%

\textbf{Calcul aproximativ de sume de serii}:

2. Calculați cu eroare $\varepsilon$ sumele seriilor:
\begin{enumerate}[(a)]
\item $x_n = \frac{(-1)^n}{n!}, \varepsilon = 10^{-3}$ \hfill (n = 7);
\item $x_n = \frac{(-1)^n}{n^3 \sqrt{n}}, \varepsilon = 10^{-2}$ \hfill (n = 4);
\end{enumerate}

\textbf{Șiruri de funcții}:

3. Studiați convergența punctuală și convergența uniformă pentru șirurile
de funcții:
\begin{enumerate}[(a)]
\item $f_n : [-1, 1] \to \RR, f_n(x) = \frac{x}{1 + nx^2}$;
\item $f_n : [0, 1] \to \RR, f_n(x) = x^n(1 - x^n)$;
\item $f_n : (-1, 1) \to \RR, f_n(x) = \frac{1 - x^n}{1 - x}$;
\item $f_n : [0, 1] \to \RR, f_n(x) = \frac{2nx}{1 + n^2 x^2}$;
\item $f_n : \RR \to \RR, f_n(x) = \arctan(nx)$.
\end{enumerate}

\vspace{1cm}

4. Verificați dacă șirul de funcții poate fi integrat termen cu termen:
\[
  f_n : [0, 1] \to \RR, \quad f_n(x) = nxe^{-nx^2}.
\]
Arătați, deci, că:
\[
  \lim_{n \to \infty} \int_0^1 f_n(x) \neq \int_0^1 \lim_{n \to \infty} f_n(x) dx.
\]

\vspace{1cm}

5. Verificați dacă șirul de funcții poate fi derivat termen cu termen:
\[
  f_n : \RR \to \RR, \quad f_n(x) = \frac{1}{n} \arctan x^n.
\]
Arătați, deci, că:
\[
  \Big( \lim_{n\to \infty} f_n(x) \Big)^\prime_{x = 1} \neq %
  \lim_{n \to \infty} f'_n(1).
\]

%%%%%%%%%%%%%%%%%%%%%%%%%%%%%%%%%%%%%%%%%%%%%%%%%%%%%%%%%%%%%%%%%%%%%%

\textbf{Serii Taylor}

6. Să se dezvolte în serie Maclaurin următoarele funcții, precizînd și
domeniul de convergență:
\begin{enumerate}[(a)]
\item $f(x) = e^x$;
\item $f(x) = \sin x$;
\item $f(x) = \cos x$;
\item $f(x) = \ln(1 + x)$;
\item $f(x) = \frac{1}{x + 1}$;
\item $f(x) = \arctan x$.
\end{enumerate}

\vspace{1cm}

7. Calculați, cu ajutorul seriilor Taylor, cu o eroare de $10^{-3}$:
\begin{enumerate}[(a)]
\item $\ds\int_0^{\frac{1}{2}} \frac{\ln ( 1 + x )}{x} dx$;
\item $\ds\int_0^{\frac{1}{3}} \frac{\sin x}{x} dx$;
\item $\ds\int_0^{\frac{1}{3}} \frac{\arctan x}{x} dx$.
\end{enumerate}

%%%%%%%%%%%%%%%%%%%%%%%%%%%%%%%%%%%%%%%%%%%%%%%%%%%%%%%%%%%%%%%%%%%%%%

\textbf{Serii de puteri}:

8. Găsiți raza de convergență și domeniul de convergență pentru seriile:
\begin{enumerate}[(a)]
\item $\sum x^n$;
\item $\sum n^n x^n$;
\item $\sum \frac{(x + 3)^n}{n^2}$.
\end{enumerate}

%%%%%%%%%%%%%%%%%%%%%%%%%%%%%%%%%%%%%%%%%%%%%%%%%%%%%%%%%%%%%%%%%%%%%%

\textbf{Serii de funcții}:

9. Studiați convergența seriilor de funcții:
\begin{enumerate}[(a)]
\item $\sum \frac{nx}{1 + n^5 x^2}, x \in \RR$; \hfill (Weierstrass);
\item $\sum \arctan \frac{2x}{x^2 + n^4}, x \in \RR$; \hfill (Weierstrass);
\item $\sum \frac{\sin nx}{2^n}, x \in \RR$; \hfill (Weierstrass);
\item $\sum \Big( \sin \frac{x}{n + 1} - \sin \frac{x}{n} \Big)$; \hfill (șirul sumelor
  parțiale);
\item $x + \sum \Big( \frac{x}{1 + nx} - \frac{x}{1 + (n-1) x} \Big), 0 \leq x \leq 1$;
  \hfill (șirul sumelor parțiale).
\end{enumerate}





\end{document}

